\documentclass[11pt,a4paper]{article}

% ─────────────────────────────────────────────────────────────────────────────
% Two claims: the algorithm as stated in Colbrook 2026 §3.4.1 cannot be followed
% by a finite element method, and what is done for Maxwell in §3.4.4 is the
% functional analytic setting proposed here instead.
% Source: notes/hand_notes.pdf.
% ─────────────────────────────────────────────────────────────────────────────

\usepackage[margin=2.2cm]{geometry}
\usepackage{amsmath,amssymb,amsthm,mathtools}
\usepackage{tikz-cd}
\usepackage[colorlinks=true,linkcolor=blue!55!black,citecolor=blue!55!black,
            urlcolor=blue!55!black]{hyperref}
\usepackage[nameinlink]{cleveref}
\usepackage{todonotes}
\newcommand{\pef}[1]{\todo[inline, color=green!40]{PEF: #1}}

\newcommand{\R}{\mathbb{R}}
\newcommand{\C}{\mathbb{C}}
\newcommand{\norm}[1]{\left\|#1\right\|}
\newcommand{\abs}[1]{\left|#1\right|}
\newcommand{\dual}[2]{\left\langle #1, #2 \right\rangle}
\newcommand{\Hb}{\mathcal{H}}
\newcommand{\ii}{\mathrm{i}}
\newcommand{\bcurl}{\bar\nabla\times}
\newcommand{\hcurl}{\hat\nabla\times}
\DeclareMathOperator{\dom}{D}
\DeclareMathOperator{\Sp}{Sp}
\DeclareMathOperator{\divg}{div}
\DeclareMathOperator{\curl}{curl}
\newcommand{\siginf}{\sigma_{\inf}}
\newcommand{\Astr}{A_{\mathrm{s}}}

% The six lettered items of section 2 are \paragraph blocks.  The default 3.25ex
% of lead before each is more than this note needs.
\makeatletter
\renewcommand\paragraph{\@startsection{paragraph}{4}{\z@}%
  {1.4ex \@plus0.4ex \@minus.2ex}{-1em}{\normalfont\normalsize\bfseries}}
\renewcommand\section{\@startsection{section}{1}{\z@}%
  {-2.5ex \@plus -1ex \@minus -.2ex}{1.6ex \@plus.2ex}%
  {\normalfont\Large\bfseries}}
\renewcommand\subsection{\@startsection{subsection}{2}{\z@}%
  {-2.2ex \@plus -1ex \@minus -.2ex}{1.2ex \@plus.2ex}%
  {\normalfont\large\bfseries}}
\makeatother

\setlength{\emergencystretch}{1em}

% Compact bibliography, so the references sit on the last page of text.
\let\oldthebibliography\thebibliography
\renewcommand{\thebibliography}[1]{%
  \small\oldthebibliography{#1}\setlength{\itemsep}{0pt}\setlength{\parskip}{0pt}}

\title{\vspace{-1.5em}The functional analytic aspects of operator folding\\
       and finite element methods\vspace{-0.5em}}
\author{Umberto Zerbinati\thanks{Mathematical Institute, University of Oxford.
  Corresponding author,
  \href{mailto:zerbinati@maths.ox.ac.uk}{zerbinati@maths.ox.ac.uk}.}
  \and Margaret McCarthy\footnotemark[1]
  \and Patrick E. Farrell\footnotemark[1]\vspace{-0.5em}}
\date{\today}

\begin{document}
\maketitle

\noindent
In \cite{Colbrook2026} Colbrook introduces very appealing algorithms for
computing spectra and pseudospectra. This note makes two points regarding its application in
the context of the finite element method.

\smallskip
\noindent\textbf{(a)}\ \emph{It is not straightforward to apply he algorithm as stated in
\cite[\S3.4.1]{Colbrook2026} in the context of a finite element method.}  The algorithm
is phrased in the language of spectral theory, in terms of e.g.~$A$, $A^*$, $\dom(A)$, $\dom(A^*)$ and orthogonal
projections $\mathcal{P}_n$. These are not the natural language in which finite element discretisations are expressed. \Cref{sec:stuck} lists some
difficulties we have encountered in passing from one setting to the other.

\smallskip
\noindent\textbf{(b)}\ \emph{What is actually computed for Maxwell in
\cite[\S3.4.4]{Colbrook2026} is not that algorithm.}  It is, we believe, the
functional analytic setting we set out in \cref{sec:path}, namely a conforming
Rayleigh--Ritz discretisation of the pencil $\eqref{eq:starstar}\,u = \lambda Mu$
on $V \times V$.  \Cref{sec:already} gives our evidence.
\pef{I can't parse this point. How can a functional analytic setting be computed for Maxwell in Colbrook's book?}

If (b) is right the two are reconciled \pef{What two? What reconciliation? Confused here}, and what is missing from
\cite{Colbrook2026} is the bridge between them, which is \cref{sec:path}.

% ═════════════════════════════════════════════════════════════════════════════
\section{The setting of Colbrook}
\label{sec:setting}

We exclusively consider the folding strategy.
Folding starts from a closed, densely defined $A$ on a Hilbert space $\Hb$.
Since we want its spectrum, $A$ must be an \emph{endomorphism}, $A \colon
\dom(A) \subset \Hb \to \Hb$, the classical instance being $\Delta \colon
H^2(\Omega) \cap H^1_0(\Omega) \to L^2(\Omega)$.
\pef{This isn't obviously an endomorphism, because $L^2$ will contain e.g.~$C^0$ functions whose boundary
trace is not zero?}
Under suitable conditions on
$\dom(A)$ and $\dom(A^*)$ one considers \cite[\S3.4.1]{Colbrook2026}
\begin{equation}
  \label{eq:gamma}
  \gamma(z,A) := \min\bigl\{\siginf(A - zI),\ \siginf\bigl((A-zI)^*\bigr)\bigr\}
    = \norm{(A-zI)^{-1}}^{-1},
\end{equation}
where the injection moduli come from the quadratic forms of the resolvent and
its adjoint,
\begin{equation}
  \label{eq:quotients}
  \siginf(A - zI)^2 = \!\!\inf_{f \in \dom(A)\setminus\{0\}}\!\!
     \frac{\norm{(A-zI)f}_{\Hb}^2}{\norm{f}_{\Hb}^2},
  \qquad
  \siginf\bigl((A-zI)^*\bigr)^2 = \!\!\inf_{f \in \dom(A^*)\setminus\{0\}}\!\!
     \frac{\norm{(A-zI)^*f}_{\Hb}^2}{\norm{f}_{\Hb}^2} .
\end{equation}
So far everything is formal and we are happy.  Our difficulty begins with the
comment below \cite[Eq.~(3.26)]{Colbrook2026}, that this reduces the computation
of pseudospectra to the eigenvalues of the self-adjoint operator
\begin{equation}
  \label{eq:folded}
  (A - zI)^*(A - zI).
\end{equation}

% ═════════════════════════════════════════════════════════════════════════════
\section{Why we cannot follow the algorithm}
\label{sec:stuck}

\paragraph{(i) The conditions on $\dom(A)$ and $\dom(A^*)$ are left unstated.}
The sentence introducing \cite[Eq.~(3.25)]{Colbrook2026} reads ``\emph{under
suitable conditions on $\dom(A)$ and $\dom(A^*)$}, the functions $\Phi_n(z,A)$
converge locally uniformly to $\gamma(z,A)$ \emph{from above} and hence satisfy
Assumption~3.2.1''. \pef{Have you checked the final version of the book? And mention what $\Phi_n$ is.} That hypothesis carries everything, since convergence from
above is what Assumption~3.2.1 asks for and what turns
\cite[Thm~3.2.3]{Colbrook2026} into $\Gamma^\epsilon_n(A) \subseteq
\Sp_\epsilon(A)$, the error control. \pef{This previous sentence strongly reads like AI.} We cannot find in the book what these
assumptions are. Reading back
from \cite[Alg.~3.4]{Colbrook2026}, which asks for vectors whose span is a core
of $A$ and of $A^*$, we take it to be
\begin{equation}
  \label{eq:suitable}
  \operatorname{ran}(\mathcal{P}^*_n) \subseteq \dom(A) \cap \dom(A^*)
  \ \text{ for each } n,
  \qquad
  \textstyle\bigcup_n \operatorname{ran}(\mathcal{P}^*_n)
  \ \text{ a core of } A \text{ and of } A^* ,
\end{equation}
the first half being needed for $\Phi_n$ to be defined at all.  On that reading
one and the same finite-dimensional space must lie in \emph{both} domains (of $A$ and $A^*$), and
those carry different boundary conditions in general.  If instead the conditions
are meant to make \eqref{eq:folded} well posed, they are different ones. \pef{What is `they'?}
\Cref{eq:folded} is the operator $T^*T$ of \cite[\S V.3.7]{Kato1995}, and by
\cite[Thm~V.3.24, p.~275]{Kato1995} it needs the strictly smaller $\dom(T^*T) =
\{u \in \dom(T) : Tu \in \dom(T^*)\}$ with $T = A - zI$, which a discretisation
would have to respect instead.

\paragraph{(ii) Babu\v{s}ka--Osborn does not apply.} 
\pef{So what? He never said Babu\v{s}ka--Osborn should apply?}
We would like to apply
Babu\v{s}ka--Osborn theory \cite{BabuskaOsborn1991,Boffi2010} to
\eqref{eq:folded}, and cannot, since it is unbounded with non-compact inverse.
We take this to be deliberate, as \cite[p.~129]{Colbrook2026} says ``it is not
necessary to assume that the operator has a compact resolvent''.  What we cannot
see is what plays its role.  Rayleigh--Ritz on a conforming subspace bounds
$\gamma$ from above for free, which we take to be the source of the error
control, but it gives no rate and no local indicator, and
\cite[Prop.~3.3.12]{Colbrook2026} gives both only for \emph{normal} $A$ with an
isolated eigenvalue and a computable gap.  But we wish to consider non-normal
operators, since we are trying to compute pseudospectra?

\paragraph{(iii) A primal-dual mismatch.}  Read \eqref{eq:folded} literally, a product of $A-zI$ with its adjoint.  The Banach adjoint is
$(A-zI)^* \colon \Hb^* \to \Hb^*$, so we would be composing $\Hb
\xrightarrow{A-zI} \Hb$ with $\Hb^* \xrightarrow{(A-zI)^*} \Hb^*$, and the two
halves do not meet. \pef{Won't his obvious reply be `I am considering the Hilbert adjoint'?}
The standard repair, when $\Hb = L^2$, is to identify $\Hb$
with its dual silently.  That does keep the composition an endomorphism. Silently identifying
a space with its dual is evil, for the reasons we now outline.

\paragraph{(iv) What the identification hides.}  It makes the Riesz map coincide
with the identity.  Writing $M \colon \Hb \to \Hb^*$ for the map it really is,
the operator ought to be
\begin{equation}
  \label{eq:star}
  M^{-1}\,(A - zI)^{*}\,M\,(A - zI),
  \tag{$*$}
\end{equation}
that is $\Hb \xrightarrow{A-zI} \Hb \xrightarrow{M} \Hb^* \xrightarrow{(A-zI)^*}
\Hb^* \xrightarrow{M^{-1}} \Hb$.
Undoing the identification therefore only exhibits the Hilbert adjoint as
$M^{-1}(\,\cdot\,)^*M$ applied to the Banach adjoint, and \eqref{eq:star} is
\eqref{eq:folded} again, since $(M^{-1}(A-zI)^*M(A-zI)u, v)_{\Hb} = \dual{M(A -
zI)u}{(A-zI)v} = ((A-zI)u, (A-zI)v)_{\Hb}$.  So we are no better off.  Finite
element operators go into $\Hb^*$, not $\Hb$, and we do not see how to read a
variational formulation off \eqref{eq:star}.

\paragraph{(v) The trial space must be a core in the graph norm.}  Points (i) to
(iv) are local.  The structural difficulty is that folding evaluates
$\norm{(A-zI)u}_{\Hb}$, so the trial space must sit inside the \emph{graph}
space $\dom(A)$ and exhaust a core of it.  For $\Delta$ on a bounded $\Omega$
four things then go wrong at once.  Continuous piecewise linears are not in
$H^2(\Omega)$, so $\norm{(A - zI)u_h}_{L^2}$ is not inaccurate but undefined, and
$C^1$ elements would be needed.  On a non-convex polygon even those fail, since
$\dom(\Delta) = \{u \in H^1_0 : \Delta u \in L^2\}$ contains corner
singularities and has $H^2 \cap H^1_0$ as a proper \emph{closed} subspace in the
graph norm, so no $H^2$-conforming family is a core.  Nor is $C^\infty_c(\Omega)$
a core, its graph closure being $H^2_0(\Omega)$, so the reflex that smooth
functions are dense is unavailable.  And $\dom(A^*)$ carries different boundary
conditions, so the second quotient in \eqref{eq:quotients} needs a second space.
None of this arises in \cite[\S\S3.4.2--3.4.3]{Colbrook2026}, set on $L^2(\R^d)$
where $C^\infty_c(\R^d)$ is \emph{assumed} to be a core and the basis is Hermite
functions, which are smooth, with no boundary.  Nor in
\cite[\S3.4.4]{Colbrook2026}, whose operator is the \emph{first-order} Maxwell
system, whose graph space $H_0(\curl) \times H(\curl)$ standard elements reach.
Between them the two settings never cover a second-order operator on a bounded
domain with boundary conditions.  The same conclusion is reached from the
discrete side in \cite[Rmk~3.7]{LazzarinoNakatsukasaZerbinati2026}, which warns
that the $L^2$ Riesz map ``will not be consistent with the function spaces of
the continuous problem''.

\paragraph{(vi) The estimator we want needs more than either.}  We intend to
drive refinement by a dual weighted residual estimator
\cite{BeckerRannacher2001,HeuvelineRannacher2001}.  DWR needs a form bounded on
$V \times V$, a conforming $V_h$ for Galerkin orthogonality, a well posed dual
problem, and, for an eigenvalue problem, an isolated eigenvalue to linearise
about.  The third is free, the folded operator being self-adjoint so that the
dual problem is the primal one.  The first two are (v) again.  The fourth is the
hard one.  With \emph{essential} spectrum at the bottom there is no
eigenfunction, hence no dual problem and nothing to linearise about.  The
discrete minimiser still exists and the upper bound of (ii) still holds, so the
certificate survives, but the estimator has no anchor.  The tool needs precisely
the isolation that \cite[p.~129]{Colbrook2026} declines to assume.  Whether the
bottom is isolated also depends on where $z$ sits, so the inner and outer layers
of adaptivity are not independent.

% ═════════════════════════════════════════════════════════════════════════════
\section{The setting we propose}
\label{sec:path}

Let $V \subset\subset \Hb$ with $\Hb = L^2$, taking $V$ of the type $H_0(\curl)
\cap H(\divg 0)$ so the embedding is compact.\footnote{Two caveats.  We are
sloppy about the components: $E$ and $H$ carry different boundary conditions, so
the pair lives in $H_0(\curl) \times H(\curl)$ rather than $V \times V$.  And the
solenoidal constraint buys the compact embedding at the cost of conformity,
since standard elements are not divergence-free, as flagged in
\cite[footnote~2, p.~130]{Colbrook2026}.}  From here $A$ is the \emph{weak}
operator, related to the endomorphism of \cref{sec:setting} by $\Astr = M^{-1}A$,
and only now does the shift $A - zM$ make sense.  Take
\begin{equation}
  \label{eq:maxwellop}
  A \colon V \times V \to V^* \times V^*,
  \qquad
  \dual{A\begin{psmallmatrix} E \\ H \end{psmallmatrix}}
       {\begin{psmallmatrix} G \\ F \end{psmallmatrix}}
    = \bigl(-\ii\,\bcurl E,\ F\bigr)_{L^2}
      + \bigl(\ii\,\hcurl H,\ G\bigr)_{L^2} .
\end{equation}
Since $A$ is \emph{first order} it goes not into $V^*$ but into $\Hb^* \subseteq
V^*$.  So we may bring in the Riesz map $M^{-1} \colon \Hb^* \to \Hb$ and write
the folded operator as
\begin{equation}
  \label{eq:starstar}
  (A - zM)^{*}\,M^{-1}\,(A - zM),
  \tag{$**$}
\end{equation}
which composes, $V \times V \xrightarrow{A - zM} \Hb^* \times \Hb^* \subseteq
V^* \times V^* \xrightarrow{M^{-1}} \Hb \times \Hb \cong (\Hb^* \times \Hb^*)^*
\xrightarrow{(A-zM)^*} V^* \times V^*$, the middle step being the canonical
double dual identification, numerically an identity.  This is the continuous
counterpart of the modified normal equation $A^{T}TAx = A^{T}Tb$ of
\cite[Eq.~(3.7)]{LazzarinoNakatsukasaZerbinati2026}, where $T$ discretises a
Riesz map inserted for the reason given in \cref{sec:stuck}(iii), that the plain
normal equation acts on inconsistent function spaces.  Folding adds that the
Riesz map is not free.  It must be that of the pivot space $\Hb$, where
$\siginf$ is measured.
Unlike \eqref{eq:star}, this has a variational form attached,
\begin{multline}
  \label{eq:varform}
  a(E,F,H,G) = \bigl(\bcurl E,\ \bcurl F\bigr)
             + \bigl(\hcurl H,\ \hcurl G\bigr)
             - (\bar z + z)\,\ii \bigl(\hcurl H,\ F\bigr) \\
             + (\bar z + z)\,\ii \bigl(\bcurl E,\ G\bigr)
             + \bar z z\,(E,F) + \bar z z\,(H,G).
\end{multline}

\subsection{The G\aa rding inequality, and the standard finite element device}
\label{sec:garding}

\Cref{sec:stuck} got stuck because \cite{Colbrook2026} needs an endomorphism
while a finite element method produces $A \colon V \to V^*$, an endomorphism of
nothing, with no spectrum of its own.  Above we repaired that by
\emph{composing} with a Riesz map, recovering the unbounded $\Astr = M^{-1}A$ on
$\Hb$.  A finite element method usually does the opposite and \emph{inverts},
and what comes out is compact rather than unbounded.  Two steps, the first
free.

A G\aa rding inequality is \emph{automatic}, with no hypothesis beyond $V$
carrying the graph norm $\norm{u}_V^2 = \norm{u}_{\Hb}^2 + \norm{\Astr
u}_{\Hb}^2$.  Applying $(s-t)^2 \ge \tfrac12 s^2 - t^2$ with $s = \norm{\Astr
u}_{\Hb}$ and $t = \abs{z}\,\norm{u}_{\Hb}$ gives, for $u \in V$,
\begin{equation}
  \label{eq:garding}
  a(u,u) = \norm{(\Astr - zI)u}_{\Hb}^2
    \ \ge\ \tfrac12 \norm{\Astr u}_{\Hb}^2 - \abs{z}^2\norm{u}_{\Hb}^2
    \ =\ \tfrac12 \norm{u}_V^2 - \bigl(\tfrac12 + \abs{z}^2\bigr)
         \norm{u}_{\Hb}^2 ,
\end{equation}
so $a + C(\cdot,\cdot)_{\Hb}$ is coercive on $V$ for $C = \tfrac12 + \abs{z}^2$,
the constant $\tfrac12$ being unimportant beyond its independence of $z$.

Lax--Milgram then makes $\eqref{eq:starstar} + CM \colon V \to V^*$ an
isomorphism, so the solution operator
\begin{equation}
  \label{eq:solop}
  T \colon \Hb \to V, \qquad
  a(Tf, v) + C\,(Tf, v)_{\Hb} = (f,v)_{\Hb}
  \quad \text{for all } v \in V,
\end{equation}
is well defined, bounded, self-adjoint and positive.  Composed with
$V \subset\subset \Hb$ it is a
\emph{compact endomorphism of $\Hb$}, and this, not $\eqref{eq:starstar}$, is
the object whose spectrum a finite element method studies.  Its nonzero
eigenvalues are reciprocal to those of the pencil, $\lambda = \mu^{-1} - C$, so
Riesz--Schauder gives a pencil spectrum $0 \le \lambda_1 \le \lambda_2 \le \cdots
\to +\infty$ of finite multiplicities, equivalently
$\eqref{eq:starstar} - \lambda M$ Fredholm of index zero for every $\lambda$.
There is no essential spectrum, the bottom $\lambda_1$ is an isolated eigenvalue,
equal to $\siginf(\Astr - zI)^2$ when $V$ is the whole graph space, and
Babu\v{s}ka--Osborn theory \cite{BabuskaOsborn1991,Boffi2010} applies to $T$ in
the usual way, bringing Kato--Temple and hence DWR with it.

The device does not solve the difficulty of \cref{sec:setting}, it inverts it.
An endomorphism is produced by \emph{solving} rather than by \emph{applying},
trading the unbounded $\Astr$ for the compact $T$ with reciprocal spectrum.  The
G\aa rding inequality is what makes $T$ exist, the compact embedding is what
makes it compact.

We know what this costs.  A compact embedding $V \subset\subset \Hb$ of the graph
space is the same thing as compactness of the resolvent of $\Astr$, exactly what
\cite[p.~129]{Colbrook2026} declines to assume.  For Maxwell it forces the
solenoidal $V$ of the footnote, which standard finite elements do not provide.
That trade seems to us the crux of the whole question.

\subsection{This appears to be what is already done for Maxwell}
\label{sec:already}

Read \cite[\S3.4.4]{Colbrook2026} back through the above.  The quantity minimised
in step~2 on \cite[p.~132]{Colbrook2026} is
\begin{equation}
  \label{eq:colbrookmin}
  \min_{(E,H) \in V_N \setminus \{0\}}
    \frac{\int_U \abs{\ii\,\hcurl H - zE}^2
                 + \abs{-\ii\,\bcurl E - zH}^2 \,\mathrm{d}x}
         {\int_U \abs{E}^2 + \abs{H}^2 \,\mathrm{d}x}
  \ =\ \min_{u_h \in V_N\setminus\{0\}}
    \frac{\norm{(\Astr - zI)u_h}^2_{L^2}}{\norm{u_h}^2_{L^2}} ,
\end{equation}
over a conforming $V_N \subset V$ of \emph{plain Lagrange elements}.  Choosing a
basis of $V_N$, the numerator is the stiffness matrix of \eqref{eq:varform},
that is of \eqref{eq:starstar}, and the denominator is the mass matrix.  So
\eqref{eq:colbrookmin} is precisely the Rayleigh--Ritz discretisation of the
pencil $\eqref{eq:starstar}\,u = \lambda M u$.

What is never built is as telling as what is.  The operator \eqref{eq:folded} is
not formed.  No adjoint is applied to anything.  No domain of Kato type is
discretised.  No projection $\mathcal{P}_n$ appears.  Only a bilinear form and a
mass matrix are assembled, which is all a finite element code can offer.  Since
the Maxwell operator is self-adjoint the second quotient in
\eqref{eq:quotients} coincides with the first and is skipped, which is why the
two-space difficulty of \cref{sec:stuck}(i) does not surface there either.

This is why we suspect the algorithm as stated in \cite[\S3.4.1]{Colbrook2026}
is not what a finite element method can follow, and that
\cite[\S3.4.4]{Colbrook2026} follows \eqref{eq:starstar} instead.  We are not
suggesting that section is wrong.  We are suggesting it does something the
statement of the algorithm does not say, and that the missing step is
\cref{sec:path}.

% ═════════════════════════════════════════════════════════════════════════════
\section{Open questions}
\label{sec:questions}

\begin{enumerate}\itemsep1pt
  \item \textbf{The main one.}  Is the computation in
        \cite[\S3.4.4]{Colbrook2026} already the reading \eqref{eq:starstar},
        rather than the algorithm of \cite[\S3.4.1]{Colbrook2026}?  If so, is the
        passage between the two the one given in \cref{sec:path}, and can the
        algorithm be restated in those terms?
  \item Which reading of \cref{sec:stuck}(i) is meant, \eqref{eq:suitable}, so
        that only the forms \eqref{eq:quotients} on $\dom(A)$ and $\dom(A^*)$ are
        ever needed, or Kato's $\dom(T^*T)$, so that \eqref{eq:folded} itself
        must be realised?
  \item Is the mismatch of \cref{sec:stuck}(iii) real, and is the Riesz map in
        \eqref{eq:starstar} the right repair?  Is it forced to be the Riesz map
        of the pivot space $\Hb$ rather than that of $V$?
  \item \Cref{sec:garding} removes the essential spectrum that blocks DWR in
        \cref{sec:stuck}(vi), at the price of the compact embedding, which is
        compactness of the resolvent and is what \cite{Colbrook2026} deliberately
        avoids.  Is that trade the right one when the goal is a finite element
        method with a computable local error indicator?  If not, what drives
        refinement when the bottom of the folded spectrum is essential?
\end{enumerate}

\bibliographystyle{plain}
\bibliography{references}

\end{document}
